% 2.11 Caso de uso - Gestión de Análisis de Germinación
\subsubsection{Gestión de Análisis de Germinación}

\noindent\textbf{Caso de uso:} Crear y Gestionar Tablas de Germinación

\vspace{0.5em}
\noindent\textbf{Descripción:} El usuario crea tablas de germinación con condiciones experimentales específicas, agrega repeticiones, registra conteos por fecha, y el sistema calcula automáticamente totales y porcentajes.

\vspace{0.5em}
\noindent\textbf{Pre-Condición:}
\begin{itemize}[nosep]
    \item El análisis de Germinación debe existir con estado \texttt{EN\_PROCESO}
    \item El usuario debe estar autenticado
    \item El usuario debe tener permisos para crear tablas
\end{itemize}

\vspace{0.5em}
\noindent\textbf{Post-Condición:}
\begin{itemize}[nosep]
    \item Tabla(s) de germinación creada(s) con configuración específica
    \item Repeticiones creadas y conteos registrados
    \item Cálculos automáticos realizados (totales, promedios, porcentajes)
    \item Análisis listo para finalizar cuando hay al menos una tabla completa
\end{itemize}

\vspace{0.5em}
\noindent\textbf{Actores:} Administrador.

\vspace{0.5em}
\noindent\textbf{Flujo de evento principal:}
\begin{enumerate}[nosep]
    \item \textbf{Fase 1: Creación de tabla de germinación}
    \begin{enumerate}[nosep]
        \item El usuario accede al detalle del análisis de Germinación
        \item Selecciona ``Crear nueva tabla de germinación''
        \item Configura parámetros de la tabla:
        \begin{itemize}[nosep]
            \item Fecha de inicio
            \item Fechas de conteos (hasta 4 fechas)
            \item Número de repeticiones: entre 2 y 8
            \item Número de semillas por repetición: 25, 50 o 100
            \item Prefrío (días, opcional)
            \item Pretratamiento (descripción, opcional)
            \item Método: sustrato, entre papel, sobre papel, etc.
            \item Temperatura (°C)
            \item Tratamientos adicionales (opcional)
        \end{itemize}
        \item Guarda la configuración
        \item El sistema crea la tabla con estado inicial
    \end{enumerate}

    \item \textbf{Fase 2: Creación de repeticiones}
    \begin{enumerate}[nosep]
        \item El usuario selecciona ``Agregar repetición'' a la tabla
        \item El sistema crea la repetición (ej. Repetición 1, 2, 3...)
        \item La repetición queda lista para registrar conteos
        \item Repite hasta alcanzar el número de repeticiones configurado
    \end{enumerate}

    \item \textbf{Fase 3: Registro de conteos}
    \begin{enumerate}[nosep]
        \item Para cada repetición, el usuario ingresa conteos por fecha:
        \begin{itemize}[nosep]
            \item Fecha 1: número de semillas germinadas normalmente
            \item Fecha 2: número de semillas germinadas normalmente
            \item Fecha 3: número de semillas germinadas normalmente (si aplica)
            \item Fecha 4: número de semillas germinadas normalmente (si aplica)
        \end{itemize}
        \item Además puede registrar:
        \begin{itemize}[nosep]
            \item Semillas anormales
            \item Semillas muertas
            \item Semillas duras
            \item Semillas frescas
        \end{itemize}
    \end{enumerate}

    \item \textbf{Fase 4: Cálculos automáticos}
    \begin{enumerate}[nosep]
        \item El sistema calcula automáticamente por repetición y conteo:
        \begin{itemize}[nosep]
            \item Total de semillas germinadas por fecha
            \item Porcentaje de germinación = (germinadas / total semillas) $\\times$ 100
            \item Totales acumulados por conteo
        \end{itemize}
        \item Calcula promedios entre repeticiones:
        \begin{itemize}[nosep]
            \item Promedio de germinación normal por fecha
            \item Porcentaje promedio final
        \end{itemize}
    \end{enumerate}

    \item \textbf{Fase 5: Gestión de múltiples tablas}
    \begin{enumerate}[nosep]
        \item El usuario puede crear múltiples tablas con diferentes condiciones experimentales
        \item Cada tabla es independiente y tiene sus propios cálculos
        \item Permite comparar resultados entre diferentes condiciones
    \end{enumerate}
\end{enumerate}

\vspace{0.5em}
\noindent\textbf{Flujos alternativos:}
\begin{itemize}[nosep]
    \item \textit{Intentar finalizar análisis sin tablas:}
    \begin{itemize}[nosep]
        \item El sistema valida que exista al menos una tabla
        \item Bloquea finalización
        \item Muestra error: ``Debe crear al menos una tabla de germinación antes de finalizar''
    \end{itemize}

    \item \textit{Tabla sin repeticiones completas:}
    \begin{itemize}[nosep]
        \item El sistema no calcula promedios
        \item Muestra advertencia: ``Esta tabla tiene repeticiones incompletas''
        \item Permite continuar pero marca la tabla como incompleta
    \end{itemize}

    \item \textit{Conteo excede número de semillas por repetición:}
    \begin{itemize}[nosep]
        \item El sistema valida en tiempo real
        \item Muestra error: ``El conteo no puede exceder [N] semillas configuradas''
        \item Bloquea guardado hasta corregir
    \end{itemize}

    \item \textit{Editar configuración de tabla existente:}
    \begin{itemize}[nosep]
        \item El usuario puede editar parámetros de la tabla
        \item Si ya hay repeticiones con datos → muestra advertencia
        \        item Requiere confirmación para continuar
    \end{itemize}

    \item \textit{Eliminar tabla:}
    \begin{itemize}[nosep]
        \item El usuario puede eliminar tablas
        \item Si tiene repeticiones con datos → requiere confirmación
        \item Elimina en cascada todas las repeticiones y conteos asociados
    \end{itemize}
\end{itemize}

\vspace{0.5em}
\noindent\textbf{Requerimientos especiales:}
\begin{itemize}[nosep]
    \item Cada tabla es independiente con su propia configuración
    \item Cálculos automáticos en tiempo real por repetición y promediados
    \item Validación de conteos no excedan semillas configuradas
    \item Permite múltiples condiciones experimentales en un mismo análisis
    \item Fechas de conteo deben ser posteriores a fecha de inicio
    \item Número de fechas de conteo flexible (1 a 4)
\end{itemize}
